% ---------------------------------------------------------------------------
%  QML Lab research project 15: Barren-Plateau Mitigation
%  Build:  latexmk -pdf main.tex   (needs qmlproject.sty, harvardqml.bst, references.bib)
% ---------------------------------------------------------------------------
\documentclass[11pt,a4paper]{article}
\usepackage{qmlproject}
\qmlheader{Project 15: Barren-Plateau Mitigation}

\hypersetup{pdftitle={QML Lab project 15: Barren-Plateau Mitigation}}

\begin{document}

\begin{center}
  {\small\textcolor{black!55}{QML Lab \textbullet\ research project 15}}\\[0.5em]
  {\LARGE\bfseries Barren-Plateau Mitigation}\\[0.6em]
  {\large Comparing strategies that keep gradients alive as circuits grow.}\\[0.6em]
  {\small Living document \textbullet\ last updated \today}
\end{center}

\begin{projectcard}
\small
\begin{tabularx}{\linewidth}{@{}>{\bfseries}l X >{\bfseries}l X@{}}
Quadrant  & \CQ\ (classical data, quantum processor) & Track & Noise and hardware \\
Difficulty & Medium & Tools & PennyLane \\
Builds on & Module 8 of the QML Lab & Workspace & \path{qml_projects/barren-plateau-mitigation/} \\
\end{tabularx}
\end{projectcard}

\paragraph{Summary.}
Module 8 showed that gradients of random deep circuits vanish exponentially with the number of qubits.
This project compares the main remedies, local cost functions, small-angle and identity-block
initialisation, and layerwise training, by gradient variance at initialisation and by what they achieve
after a fixed training budget. It also asks the uncomfortable question raised by recent theory: are the
circuits that avoid barren plateaus the ones a classical computer can simulate?

% ===========================================================================
\section{Background and key concepts}
% ===========================================================================
\subsection{Barren plateaus}
For a cost $C(\theta)=\Tr[O\,U(\theta)\rho U(\theta)^\dagger]$ and sufficiently random circuits,
\citet{mcclean2018} showed
\begin{equation}
  \E_\theta[\partial_kC]=0,\qquad \operatorname{Var}_\theta[\partial_kC]\in O(b^{-n}),\ b>1 .
\end{equation}
Exponentially many shots would be needed to estimate a gradient. \citet{cerezo2021cost} showed that the
choice of cost matters: global observables (acting on all qubits) give barren plateaus even for shallow
circuits, local observables on shallow circuits keep gradients polynomially large. \citet{holmes2022}
linked the effect to expressibility: the closer the circuit ensemble is to random unitaries, the flatter the
landscape. Noise adds its own plateau \citep{wang2021noise}.

\subsection{Mitigation strategies}
\begin{description}[style=nextline,leftmargin=1.2em]
  \item[Identity-block initialisation] \citet{grant2019} build the circuit from blocks
        $U_b(\theta_b)\,U_b(\theta_b')^\dagger$ initialised with $\theta_b'=\theta_b$, so the whole circuit
        starts as the identity and early gradients are large.
  \item[Small-angle initialisation] Drawing $\theta\sim\mathcal N(0,\sigma^2)$ with small $\sigma$ keeps the
        circuit near the identity at the start.
  \item[Layerwise learning] \citet{skolik2021} train a few layers at a time and freeze them, adding depth
        gradually.
  \item[Local costs] Replace a global observable such as $Z^{\otimes n}$ by a sum of local ones.
\end{description}

\subsection{The simulability caveat}
\citet{cerezo2023simulability} argue that many architectures with provable absence of barren plateaus
can be simulated efficiently classically, because their relevant dynamics lives in a small subspace. A
complete study should therefore ask, for each successful strategy, whether a classical surrogate could
reproduce it.

% ===========================================================================
\section{Research questions}
% ===========================================================================
The project is open-ended: these questions are the starting point, not the finish line. Add new
ones to the README and the journal as they come up.

\begin{rqlist}
  \item Which strategy keeps gradient variance flattest as the number of qubits grows from 2 to 12?
  \item Does a larger initial gradient translate into a lower final loss after a fixed training budget?
  \item How do strategies combine, for example local cost with identity-block initialisation?
\end{rqlist}

% ===========================================================================
\section{Method}
% ===========================================================================
\subsection{Data}
Two training tasks: learning to return a randomly prepared state to $\ket{0\cdots0}$ (a controllable
state-learning problem), and the Module 5 classifier on two moons with increasing qubit numbers.

\subsection{Models}
The hardware-efficient ansatz of Module 8 with $L=n$ layers, extended with the parameters
\texttt{init\_strategy} (uniform, small-angle, identity-block), \texttt{cost} (global, local) and
\texttt{training} (all layers, layerwise). Gradient variances from 200 random initialisations per setting.

\subsection{Baselines}
Module 8's configuration (uniform initialisation, global cost) as the reference; analytic
$O(2^{-n})$ scaling as a guide.

\subsection{Metrics}
\begin{itemize}
  \item $\operatorname{Var}[\partial_kC]$ versus $n$ on a log scale, with fitted decay rate.
  \item Final loss and steps to reach a threshold after a fixed budget.
  \item Evidence for or against classical simulability of the successful settings.
\end{itemize}

% ===========================================================================
\section{Planned experiments}
% ===========================================================================
Each experiment becomes one \texttt{[[experiments]]} entry in the workspace, so its runs, parameters
and metrics are stored and compared in the portal.

\begin{table}[h]
\centering\small
\begin{tabularx}{\linewidth}{@{}l>{\raggedright}X>{\raggedright}X>{\raggedright\arraybackslash}X@{}}
\toprule
ID & Experiment & What varies & Main metrics \\
\midrule
  E1 & Reproduce Module 8 & $n=2..12$, uniform, global & variance slope \\
  E2 & Initialisation & uniform, small-angle, identity-block & variance, final loss \\
  E3 & Cost locality & global vs local & variance, final loss \\
  E4 & Layerwise learning & layers added per stage & final loss \\
  E5 & Combinations and noise & best pairs, depolarizing $p$ & final loss \\
\bottomrule
\end{tabularx}
\caption{Experiment plan. Results go into the journal and the results log (\cref{sec:log}).}
\end{table}

% ===========================================================================
\section{Running it in the QML Lab}
% ===========================================================================
The workspace lives in \path{qml_projects/barren-plateau-mitigation/}. Scripts import three helpers from
\path{qml_lab/plots.py}:
\begin{itemize}
  \item \texttt{params()} reads the parameters from the portal's run form (defaults when run by hand);
  \item \texttt{save()} stores figures with the run;
  \item \texttt{log\_metrics()} records numbers that appear in the run table.
\end{itemize}

\begin{lstlisting}
qml_projects/barren-plateau-mitigation/
  project.toml          metadata + experiments
  README.md             goal, open questions, ideas
  experiments/          runnable scripts
  docs/main.tex         this document
\end{lstlisting}

First experiment: This experiment is \emph{planned}: write the script, add the block below to \texttt{project.toml}, then run \texttt{python manage.py sync\_projects}.

\begin{lstlisting}[language={}]
[[experiments]]
slug = "init-strategies"
title = "Gradient variance by initialisation strategy"
script = "experiments/e01_init_strategies.py"
figure = "variance_by_strategy.png"

[experiments.params]
max_qubits = 12
init_strategy = "identity_block"
cost = "local"
samples = 200
seed = 0
\end{lstlisting}

\begin{lstlisting}[language=bash]
# from the portal folder: pick up changes to project.toml
python manage.py sync_projects
cd qml_projects/barren-plateau-mitigation
# run by hand with the default parameters
PYTHONPATH=../../qml_lab python experiments/e01_init_strategies.py
\end{lstlisting}

% ===========================================================================
\section{Quality bar}
% ===========================================================================
\begin{itemize}
  \item A classical baseline with a similar number of trainable parameters, tuned with the same care \citep{bowles2024}.
  \item At least five random seeds per configuration, reported as mean $\pm$ standard deviation.
  \item How results change with the number of qubits.
  \item At least one run under a realistic noise model.
  \item Report both gradient variance at initialisation and final loss; one without the other is misleading.
\end{itemize}

% ===========================================================================
\section{Open directions}
% ===========================================================================
\begin{itemize}
  \item Symmetry-respecting (equivariant) ans\"atze.
  \item Warm starts from classical solutions or from smaller circuits.
  \item Classical surrogates for the trainable regimes, testing the simulability argument directly.
\end{itemize}

% ===========================================================================
\section{Risks and limitations}
% ===========================================================================
\begin{itemize}
  \item Large initial gradients do not guarantee good final solutions.
  \item Identity-block initialisation helps only at the start of training.
  \item State-vector memory limits the study to about 14 qubits.
\end{itemize}

% ===========================================================================
\section{Results log}
\label{sec:log}
% ===========================================================================
Summarise findings here as they become stable; day-to-day notes belong in the portal journal.

\begin{table}[h]
\centering\small
\begin{tabularx}{\linewidth}{@{}l l >{\raggedright}X >{\raggedright\arraybackslash}X@{}}
\toprule
Date & Run & Finding & Next step \\
\midrule
 & & & \\
 & & & \\
 & & & \\
\bottomrule
\end{tabularx}
\end{table}

\bibliographystyle{harvardqml}
\bibliography{references}

\end{document}
